\scanpage{144}
désignera, pour un élément
\[
(u_{-1},u_1)\in\mathbb R^{\unsure{2}}\times\mathbb R
\]
par
\[
(13)\qquad s\longmapsto u_s
\]
l'application affine de $\mathbb R$ dans $\mathbb R$ \deleted{\unsure{def}} qui prend la valeur $u_{-1}$ pour $s=-1$, la valeur $u_1$ pour $s=+1$, \deleted{\unsure{et}} donc
\[
(14)\qquad u_s=\frac{u_1-u_{-1}}2\,s+\frac{u_1+u_{-1}}2,
\]
et on pose
\[
(15)\qquad\left\{\begin{aligned}
(u_{-1},u_1).(t,s)&\overset{\text{déf}}{=}(t+u_s,s)\\
\sigma.(t,s)&\overset{\text{déf}}{=}(-t,s)
\end{aligned}\right.
\]
d'où une opération de $D_{\mathbb R\times\mathbb R}$ sur $\widetilde X$, compatible d'ailleurs avec la \deleted{\unsure{pri}} projection
\diagram{\begin{tikzcd}
\widetilde X=\mathbb R\times\unsure{\mathbb I|\mathbb J} \arrow[r] & \unsure{\mathbb I|\mathbb J}\\
(t,s) \arrow[r,mapsto] & s
\end{tikzcd}}
et l'opération triviale sur $\unsure{\mathbb I|\mathbb J}$. L'opération passe bien au quotient en une opération de $D_{\mathbb R\times\mathbb R}$ sur $X$ lui-$\overline{\mathrm m}$, compatible avec la projection
\diagram{\begin{tikzcd}
X=\mathbb U\times\unsure{\mathbb I|\mathbb J} \arrow[r] & \unsure{\mathbb I|\mathbb J}\\
(z,\unsure{s}) \arrow[r,mapsto,no head] & s
\end{tikzcd}}
et l'opération triviale sur $\unsure{\mathbb I|\mathbb J}$.

Considérons le sous-groupe de $D_{\mathbb R\times\mathbb R}$ qui opère \unsure{trivialemt|trivialement} sur $\partial X=\mathbb U\times\{-1,+1\}$, c'est le sous-groupe $\mathbb Z\times\mathbb Z\subset D_{\mathbb R\times\mathbb R}^{0}$.
\margin{144}
